% =============================================================
% Section 3. Triangle Equivalence: FEEC Geometry, Hodge Spectrum,
% and Robinson Sheaf Cohomology
% =============================================================

% --- Opening hook (continuing from Section 2) ---
We now construct the bridge promised in \cref{sec:intro}: three
formulations of well-conditioning for $k$-space sampling, drawn from
three schools that have not previously been combined---finite element
exterior calculus (FEEC) singular value bounds, topological signal
processing (TopSP) Riesz frame conditions, and Robinson sheaf
cohomology with sheaf bandwidth---are formally equivalent under the
common finite-dimensional sampling-operator normal form introduced
in \cref{subsec:tri-setup}. The equivalence is asserted in this
finite-dimensional normal form; the original infinite-dimensional
or qualitative formulations of the three source literatures are not
themselves claimed to be equivalent outside this normalization (see
\cref{subsec:tri-edge}, (L1)). The equivalence is stated as
\cref{thm:triangle-equiv} and proved in
\cref{subsec:tri-ab,subsec:tri-ac}. The proof of $(a) \Leftrightarrow
(c)$ produces a candidate definition of sheaf bandwidth that is not
categorically intrinsic; the categorically intrinsic version is
constructed in \cref{sec:sheafbw}. We close the section by formalizing
Robinson's open problem from \cite[\S 6]{Robinson2014} as a precisely
narrowed two-clause question (\cref{subsec:tri-narrowing}).

% =============================================================
\subsection{Setup and notation}\label{subsec:tri-setup}
% =============================================================

The well-conditioning of $k$-space sampling admits three equivalent
readings: an SVD-style ratio $\sigma_{\min}(A)/\sigma_{\max}(A)$ of the
sampling operator; a Riesz frame inequality on a bandlimited subspace;
and a sheaf-cohomological condition on an abstract sampling cell complex.
\Cref{thm:triangle-equiv} makes this precise: the three readings agree
on the same threshold $\tau \in (0,1]$. The setup below introduces the
common notation in which all three formulations live---the bandlimited
subspace $V_B$, the sampling operator $A$, and the condition number
$\kappa(A)$---before the triangle statement in
\cref{subsec:tri-theorem}.

Let $\Omega_k := B(0, k_{\max}) \subset \mathbb{R}^d$ be a closed ball
in $k$-space, and let $V = \{k_1, \ldots, k_N\} \subset \Omega_k$ be a
finite set of $N$ candidate sample positions. We do not impose
geometric structure on $V$ at this stage; the dependence of the
constructions on $V$ is recorded explicitly via $V_B(V)$ below and
discussed in \cref{sec:three_instances}.

Let $X(V)$ denote the Delaunay simplicial complex on $V$ (cells:
vertices $V$, edges $E$, faces $F$, tetrahedra $T$), oriented in the
sense of \cite[\S 2.1]{Hatcher2002}. Edge weights are
$w_{ij} := \lvert k_i - k_j \rvert^{-2}$ (with the convention
$w_{ij} = 0$ if $\{i,j\} \notin E$), and
$W := (w_{ij})_{1 \le i, j \le N}$. Set
$D := \mathrm{diag}(d_1, \ldots, d_N)$ with $d_i := \sum_j w_{ij}$,
and define the graph Laplacian
$L_0 := D - W \in \mathbb{R}^{N \times N}$. The eigenpairs of $L_0$
are denoted $(\mu_l, v_l)_{l = 0, \ldots, N-1}$ with
$0 = \mu_0 \le \mu_1 \le \cdots \le \mu_{N-1}$ and
$\{v_l\}_{l=0}^{N-1}$ an orthonormal basis of $\C^N$.

For a fixed bandwidth threshold $B \in (0, \mu_{N-1})$, define the
$B$-bandlimited subspace
\begin{equation}\label{eq:VB-def}
  V_B(V) := \mathrm{span}\{\, v_l : \mu_l \le B \,\} \;\subset\; \C^N,
  \qquad K := \dim V_B(V).
\end{equation}
We write $V_B$ for $V_B(V)$ when $V$ is fixed and emphasize the
$V$-dependence only when comparing across instances
(\cref{sec:three_instances}).

Fix a sampling subset $V_s \subset \{1, \ldots, N\}$ with
$M := \lvert V_s \rvert \ge K$. Let
$\Sigma_{V_s} : \C^N \to \C^M$ be the index restriction operator
$(\Sigma_{V_s} c)_n := c_{i_n}$ for $V_s = \{i_1, \ldots, i_M\}$, and
let $\iota_B : V_B \hookrightarrow \C^N$ be the canonical isometric
embedding. The \emph{sampling operator}
\begin{equation}\label{eq:A-def}
  A \;:=\; \Sigma_{V_s} \circ \iota_B \;:\; V_B \longrightarrow \C^M
\end{equation}
is the central object of this paper. In the basis
$\{v_l\}_{l=0}^{K-1}$ of $V_B$ and the canonical basis of $\C^M$, $A$
has the matrix representation
\[
  A \;=\; \bigl[\, v_l(i_n) \,\bigr]_{n = 1, \ldots, M;\, l = 0, \ldots, K-1}
  \;\in\; \C^{M \times K}.
\]

The singular values of $A$ are denoted
\[
  \sigma_{\max}(A) \;:=\; \sigma_0(A)
  \;\ge\; \sigma_1(A) \;\ge\; \cdots \;\ge\; \sigma_{K-1}(A)
  \;=:\; \sigma_{\min}(A) \;\ge\; 0.
\]
The condition number is
$\kappa(A) := \sigma_{\max}(A) / \sigma_{\min}(A) \in [1, \infty]$,
with $\kappa(A) = \infty$ when $\sigma_{\min}(A) = 0$ (rank
deficiency). Throughout, $\tau \in (0, 1]$ denotes a fixed
well-conditioning threshold; $A$ is \emph{$\tau$-well-conditioned}
when $\kappa(A) \le 1/\tau$, equivalently
$\sigma_{\min}(A) \ge \tau \, \sigma_{\max}(A)$.

% =============================================================
\subsection{Triangle equivalence theorem}\label{subsec:tri-theorem}
% =============================================================

\begin{theorem}[Triangle Equivalence]\label{thm:triangle-equiv}\footnote{Formalized in Lean 4 as \texttt{triangle\_equivalence} in \texttt{lean/MRI/Cellular/TriangleEquivalence.lean}. The Lean statement formalizes the finite-dimensional operator/sheaf core of the theorem, with $V_B$ represented by an abstract Hilbert space and $A$ by an evaluation operator; the construction of $V_B(V)$ from the Delaunay graph is treated as setup data rather than separately formalized geometry. Full correspondence: \texttt{lean/CROSSREF.md}.}
Fix $V$, $V_s$, $B$ (and hence $V_B = V_B(V)$ and $A$ as in
\cref{subsec:tri-setup}), and fix $\tau \in (0, 1]$. The following
three statements are equivalent.
\begin{enumerate}
  \item[\textup{(a)}] \emph{(FEEC singular value form.)}
        $\sigma_{\min}(A) \ge \tau \, \sigma_{\max}(A)$.
  \item[\textup{(b)}] \emph{(Hodge-spectrum Riesz form.)} There exist
        constants $0 < \alpha \le \beta$ with $\beta / \alpha \le 1/\tau$
        such that for every $c \in V_B \setminus \{0\}$,
        \begin{equation}\label{eq:riesz}
          \alpha \,\lVert c \rVert_{\C^N}
          \;\le\;
          \lVert \Sigma_{V_s} c \rVert_{\C^M}
          \;\le\;
          \beta \,\lVert c \rVert_{\C^N}.
        \end{equation}
        That is, $V_B$ forms a Riesz sequence in $\ell^2(V_s)$ with
        frame bound ratio $\beta / \alpha \le 1/\tau$.
  \item[\textup{(c)}] \emph{(Robinson sheaf form.)} On the
        \emph{abstract sampling sheaf} $\samps$ over the abstract
        sampling cell complex $X^A$, constructed in
        \cref{def:Xsamp,def:sampling-sheaf} of \cref{sec:sheafbw},
        \begin{equation}\label{eq:robinson-sheaf-form}
          H^0\!\left( X^A;\, \samps \right) \;=\; 0
          \qquad \text{and} \qquad
          \sheafBW^{\mathrm{cand}}\!\left( \samps \right)
          \;\ge\; \tau,
        \end{equation}
        where $\sheafBW^{\mathrm{cand}}$ is the candidate sheaf
        bandwidth of \cref{def:sheafbw-candidate} below.
\end{enumerate}
\end{theorem}

\begin{remark}[Three readings of one inequality]\label{rem:three-readings}
The three forms in \cref{thm:triangle-equiv} read the same scalar
inequality in three vocabularies developed in disjoint literatures:
the FEEC singular value form is the standard reading via singular
value bounds \cite{AFW2010}; the Riesz form is the TopSP reading via
$\ell^2$ frame bounds \cite{BarbarossaSardellitti2020}; the Robinson
form translates the inequality into the language of cellular sheaves
on $X(V)$ in the sense of \cite{Robinson2014}, with the sheaf
bandwidth furnishing the quantitative bound.
\end{remark}

\Cref{thm:triangle-equiv} is proved in two parts:
$(a) \Leftrightarrow (b)$ in \cref{subsec:tri-ab}, and
$(a) \Leftrightarrow (c)$ in \cref{subsec:tri-ac}. The latter relies
on the abstract sampling sheaf $\samps$ over the cell complex $X^A$,
whose explicit construction is given in
\cref{subsec:sheafBW-Xsamp} of \cref{sec:sheafbw}; the candidate
sheaf bandwidth $\sheafBW^{\mathrm{cand}}$ used in clause~\textup{(c)}
is recorded in \cref{def:sheafbw-candidate} below in terms of
$\samps$.

% =============================================================
\subsection{Proof of $(a) \Leftrightarrow (b)$}\label{subsec:tri-ab}
% =============================================================

The equivalence $(a) \Leftrightarrow (b)$ is a direct consequence of
the variational characterization of singular values together with the
isometry of $\iota_B$.

\begin{proof}[Proof of $(a) \Leftrightarrow (b)$]
Since $\iota_B : V_B \hookrightarrow \C^N$ is an isometric embedding,
$\lVert \iota_B c \rVert_{\C^N} = \lVert c \rVert_{V_B}$ for every
$c \in V_B$, and consequently
\begin{equation}\label{eq:Avariational}
  \lVert A c \rVert_{\C^M}
  \;=\;
  \lVert \Sigma_{V_s} \iota_B c \rVert_{\C^M}
  \;=\;
  \lVert \Sigma_{V_s} c \rVert_{\C^M}
  \qquad \text{for every } c \in V_B,
\end{equation}
where in the last equality $V_B$ is identified with its image
$\iota_B(V_B) \subset \C^N$. The singular values of $A$ are then
characterized variationally by
\begin{align}
  \sigma_{\max}(A)
    &\;=\; \max_{\substack{c \in V_B \\ \lVert c \rVert = 1}}
        \lVert \Sigma_{V_s} c \rVert_{\C^M},
    \label{eq:sigma-max-var} \\
  \sigma_{\min}(A)
    &\;=\; \min_{\substack{c \in V_B \\ \lVert c \rVert = 1}}
        \lVert \Sigma_{V_s} c \rVert_{\C^M}.
    \label{eq:sigma-min-var}
\end{align}

\paragraph{$(b) \Rightarrow (a)$.}
Given $\alpha, \beta$ with the bounds \cref{eq:riesz}, take infimum
and supremum over $c \in V_B$ with $\lVert c \rVert = 1$. From
\cref{eq:sigma-max-var,eq:sigma-min-var} this yields
$\sigma_{\max}(A) \le \beta$ and $\sigma_{\min}(A) \ge \alpha$. The
hypothesis $\beta / \alpha \le 1/\tau$ then gives
$\sigma_{\min}(A) / \sigma_{\max}(A) \ge \alpha / \beta \ge \tau$,
which is $(a)$.

\paragraph{$(a) \Rightarrow (b)$.}
Set $\alpha := \sigma_{\min}(A)$ and $\beta := \sigma_{\max}(A)$.
Then $\beta / \alpha = \kappa(A) \le 1/\tau$ by hypothesis.
\Cref{eq:sigma-max-var,eq:sigma-min-var} together with
\cref{eq:Avariational} yield, after homogeneity in $c$, the bound
\cref{eq:riesz} with these constants.
\end{proof}

% =============================================================
\subsection{Proof of $(a) \Leftrightarrow (c)$}\label{subsec:tri-ac}
% =============================================================

The equivalence $(a) \Leftrightarrow (c)$ is realized through the
abstract sampling sheaf $\samps$ on the cell complex $X^A$, whose
explicit construction we forward-reference to
\cref{def:Xsamp,def:sampling-sheaf} in
\cref{subsec:sheafBW-Xsamp}. Two facts about $\samps$ proved there
suffice for the present argument:
\begin{enumerate}[label=\textup{(F\arabic*)},leftmargin=*]
  \item \emph{Cochain spaces.} $C^0(\samps) \cong V_B$ (canonical
        isometry, \cref{lem:canonical-isometry}) and
        $C^1(\samps) = \C^M$
        (\cref{lem:samps-cochain-spaces}).
  \item \emph{Coboundary equals sampling operator.} The $0$-th
        coboundary $\delta^0 : C^0(\samps) \to C^1(\samps)$ of
        $\samps$ coincides with $A$ after the canonical
        identification of \cref{lem:canonical-isometry,lem:samps-cochain-spaces},
        i.e.\ $\delta^0 = A$ as linear maps $V_B \to \C^M$
        (\cref{prop:delta-equals-A}).
\end{enumerate}
Both facts depend only on the data of \cref{def:Xsamp,def:sampling-sheaf}
and on the sampling operator $A$ recorded in \cref{subsec:tri-setup};
neither depends on \cref{thm:triangle-equiv}, so the forward
references are not circular.

\begin{definition}\label{def:sheafbw-candidate}
The \emph{candidate sheaf bandwidth} of $\samps$ is
\begin{equation}\label{eq:sheafbw-candidate}
  \sheafBW^{\mathrm{cand}}(\samps)
  \;:=\;
  \frac{\sigma_{\min}(A)}{\sigma_{\max}(A)}
  \;\in\;
  [0, 1],
\end{equation}
with the convention
$\sheafBW^{\mathrm{cand}}(\samps) := 0$ when
$\sigma_{\min}(A) = 0$. This candidate references $\samps$ through
the operator $A$ derived from the chosen embedding
$V_B \hookrightarrow \C^N$ and the basis $\{v_l\}_{l=0}^{K-1}$; it is
therefore not categorically intrinsic to the cellular sheaf data of
$\samps$. The categorically intrinsic version is the spectral-gap
ratio
$\sheafBW := \sqrt{\gamma_+(\delta^0) / \Gamma_+(\delta^0)}$
constructed in \cref{sec:sheafbw}, which on $\samps$ satisfies a
strict identity with $\sigma_{\min}(A) / \sigma_{\max}(A)$ via
\cref{thm:strict-equality}.
\end{definition}

\begin{proof}[Proof of $(a) \Leftrightarrow (c)$]
By \textup{(F1)}, the coboundary $\delta^0$ of $\samps$ is well-typed
as a linear map $V_B \to \C^M$, so the identification with $A$ in
\textup{(F2)} is type-compatible. Combining \textup{(F1)} and
\textup{(F2)} gives $\delta^0 = A$ as maps $V_B \to \C^M$, hence
$H^0(\samps) = \ker \delta^0 = \ker A$, and so
$H^0(\samps) = 0$ if and only if $\sigma_{\min}(A) > 0$. By
\cref{def:sheafbw-candidate},
$\sheafBW^{\mathrm{cand}}(\samps) \ge \tau$ if and only if
$\sigma_{\min}(A) \ge \tau \, \sigma_{\max}(A)$. The conjunction of
the two conditions in \cref{eq:robinson-sheaf-form} is therefore
equivalent to $\sigma_{\min}(A) \ge \tau \, \sigma_{\max}(A)$, i.e.\
$(a)$.
\end{proof}

% =============================================================
\subsection{Reading the candidate bandwidth}\label{subsec:tri-bw-candidate}
% =============================================================

The candidate definition \cref{eq:sheafbw-candidate} is a numerical
quantity attached to the sheaf $\samps$ via the operator $A$. Three
observations describe its character.

\paragraph{(i) Range.}
$\sheafBW^{\mathrm{cand}}(\samps) \in [0, 1]$, with the value $1$
attained when $A$ is a positive scalar multiple of an isometry on
$V_B$ (the maximally well-conditioned case;
cf.\ \cref{thm:strict-equality}) and the value $0$ attained when $A$
is rank-deficient (failure of $H^0$ vanishing).

\paragraph{(ii) Three readings, one inequality.}
\Cref{rem:three-readings} identifies the three vocabularies in which
$\sigma_{\min}(A) / \sigma_{\max}(A) \ge \tau$ is read; the candidate
definition \cref{eq:sheafbw-candidate} translates the FEEC reading
verbatim into the sheaf-cohomological reading. The TopSP reading via
Riesz frame bounds is recovered through clause $(b)$ of
\cref{thm:triangle-equiv}.

\paragraph{(iii) Lack of categorical intrinsicness.}
\Cref{def:sheafbw-candidate} uses the basis $\{v_l\}_{l=0}^{K-1}$ of
$V_B$ implicit in the matrix form of $A$, i.e.\ data of the
representation $V_B \hookrightarrow \C^N$ rather than data intrinsic
to the cellular sheaf $\samps$. The categorically intrinsic
alternative, defined in \cref{sec:sheafbw} as a spectral gap ratio
of the sheaf coboundary operator, refers only to $\mathscr{A}$, not
to an ambient embedding or to a chosen basis of stalks.

% =============================================================
\subsection{Robinson's open problem: precise narrowing}\label{subsec:tri-narrowing}
% =============================================================

Robinson \cite[\S 6]{Robinson2014} writes, in the context of cellular
sheaves whose stalks are Hilbert spaces and whose cochain complexes
admit Hodge decompositions, that ``there may be a way to discuss the
concept of \emph{bandwidth} for general sheaves of Hilbert spaces.''
\Cref{thm:triangle-equiv} together with \cref{def:sheafbw-candidate}
narrows this to a two-clause open problem.

\begin{enumerate}
  \item[\textup{(R1)}] \emph{(Categorical intrinsicness.)} Define a
        sheaf bandwidth $\sheafBW(\mathscr{A})$ for cellular sheaves
        $\mathscr{A}$ of Hilbert spaces using only the cochain-complex
        data of $\mathscr{A}$ (stalks, restriction maps, coboundary),
        without reference to an ambient embedding
        $V_B \hookrightarrow \C^N$ or to a chosen basis of stalks.
  \item[\textup{(R2)}] \emph{(Reduction on the sampling instance.)}
        Whatever intrinsic definition is given in \textup{(R1)}, it
        must reduce to $\sigma_{\min}(A) / \sigma_{\max}(A)$ on the
        sampling sheaf instance, so that the candidate
        \cref{eq:sheafbw-candidate} is recovered as a special case.
\end{enumerate}

The narrowing replaces the qualitative ``find an equivalent
definition'' in \cite[\S 6]{Robinson2014} with a concrete two-clause
specification. \Cref{sec:sheafbw} answers \textup{(R1)} by the
spectral-gap ratio definition
$\sheafBW := \sqrt{\gamma_+(\delta^0) / \Gamma_+(\delta^0)}$, and
answers \textup{(R2)} by \cref{thm:strict-equality}, which establishes
the strict identity
$\sheafBW(\samps) = \sigma_{\min}(A) / \sigma_{\max}(A)$ on the
abstract sampling cell complex of \cref{sec:sheafbw}.

% =============================================================
\subsection{Edge of soundness}\label{subsec:tri-edge}
% =============================================================

\Cref{thm:triangle-equiv} is stated and proved within the
finite-dimensional setting of \cref{subsec:tri-setup}: $V$ is a
finite vertex set in $\Omega_k$, $V_B \subset \C^N$ is finite
dimensional, and $A$ acts between finite-dimensional Hilbert spaces.
Four boundary conditions delimit the scope of the theorem; each is
recorded as a project-internal limitation, fully resolved within
this paper or instantiated in a subsequent section.

\paragraph{(L1) Infinite-dimensional stalks.}
Robinson's open problem in \cite[\S 6]{Robinson2014} is stated for
sheaves of Hilbert spaces with potentially infinite-dimensional
stalks, in the setting of Hilbert complexes \cite{BruningLesch1992}.
The triangle equivalence is established here under the
finite-dimensional restriction $\dim V_B = K < \infty$. The
categorically intrinsic version of \cref{sec:sheafbw} is constructed
in the same finite-dimensional restriction, in which closedness of
cochain ranges is automatic
(see \cref{thm:closed-hilbert-complex}). The infinite-dimensional
extension is outside the scope of this paper.

\paragraph{(L2) Categorical intrinsicness of the bandwidth.}
\Cref{def:sheafbw-candidate} is not categorically intrinsic to the
sheaf $\samps$: the basis $\{v_l\}_{l=0}^{K-1}$ underlying the
matrix form of $A$ is not part of the cellular sheaf data. The
categorically intrinsic version is given in \cref{sec:sheafbw}; it
answers clause \textup{(R1)} of \cref{subsec:tri-narrowing} and is
shown to satisfy clause \textup{(R2)} on $\samps$ via
\cref{thm:strict-equality}.

\paragraph{(L3) Spectral gap as data, not hypothesis.}
The bandlimited subspace
$V_B(V) = \mathrm{span}\{v_l : \mu_l \le B\}$ presupposes a separating
gap in the spectrum of $L_0$ at $B$; otherwise the cutoff is not
stable under perturbations of $V$. In the finite-dimensional finite
vertex setting, the spectral gap is a property of the data $(V, B)$
rather than a hypothesis of \cref{thm:triangle-equiv}: the theorem
is stated for fixed $V, V_s, B$ and makes no genericity assumption
on the gap. The three instances of \cref{sec:three_instances}
instantiate the spectral-gap data concretely; the resulting
trade-off between sampling and integration is recorded in
\cref{sec:three_instances} as Empirical Proposition~8b.

\paragraph{(L4) Vertex geometry not made explicit.}
The vertex geometry of $V$ enters \cref{thm:triangle-equiv} only
through the dependence of $V_B(V)$ on $V$ (and through it the matrix
$A = [v_l(i_n)]$). The geometry is not an independent variable of
the theorem statement. The three instances of
\cref{sec:three_instances} instantiate the abstract setting in three
domain-specific geometries (simplicial complexes for TopSP,
finite-element meshes for FEM-EM, Brillouin-zone tetrahedra for
condensed matter), and Empirical Proposition~8b records the Pareto
trade-off between sampling and integration across geometries within
each instance.

% --- Closing hook (to Section 4) ---
\Cref{thm:triangle-equiv} establishes the bridge as a formal
equivalence and narrows Robinson's open problem to clauses
\textup{(R1)} and \textup{(R2)} of \cref{subsec:tri-narrowing}. The
remaining work---constructing the categorically intrinsic sheaf
bandwidth that answers \textup{(R1)}---requires the FEEC machinery,
which we review in \cref{sec:feec_review} before resuming the
construction in \cref{sec:sheafbw}.
